\section{Torus displacements, Fourier detection and binary transfer}
\label{sec:torus-fourier}

The fixed construction separates every phase vector from its half-turn
translate. We now compare this displacement with whole-row shifts,
tiny rotations, and changes of a single leaf. The Fourier argument
will use the cutoff
\[
 Q=k^{12k}.
\]
The ratio $P_v/T\leq k^{-4k}$ controls the cost of whole-row shifts,
and $R/Q=k^{-4k}$ controls the cost of tiny rotations. After proving
the comparison for an arbitrary finite-range map on the torus, we
apply the shared binary code.

\subsection{Pointwise and averaged edit costs}

For a node $v$, let $e_v\in\R^{m^k}$ have entry one on the leaves below
$v$ and zero elsewhere. At an internal node, define
\begin{equation}\label{eq:cycle-shift}
 h_v(a)=\frac a{P_v}\sum_i c_{vi}e_{vi}\pmod1
 \qquad(a\in\Z/P_v\Z).
\end{equation}
This is well defined because each $c_{vi}$ is an integer. Let $q$ be
the depth of $v$.

\begin{lemma}[Inexpensive torus translations]\label{lem:torus-costs}
For the realization fixed after Proposition~\ref{prop:good-probability},
the following estimates hold, with $q$ denoting the depth of $v$ in
the first two statements.
\begin{enumerate}[label=(\roman*)]
\item If $v$ is internal and $a$ is an integer with $0\leq a<P_v$,
then for every $w$,
\begin{equation}\label{eq:cycle-cost}
 \ED(X(w),X(w+h_v(a)))\leq2Nm^{-q}\frac{P_v}{T}.
\end{equation}
\item If $v$ is any node and $0\leq\eta\leq1/Q$, then
\begin{equation}\label{eq:tiny-cost}
 \E_w\ED(X(w),X(w+\eta e_v))\leq Nm^{-q}\frac RQ,
\end{equation}
where $w$ has normalized Haar measure on $\T^{m^k}$.
\item If $u$ is a leaf, then for every $w$ and $\eta\in\T$,
\begin{equation}\label{eq:leaf-cost}
 \ED(X(w),X(w+\eta e_u))\leq Nm^{-k}.
\end{equation}
\end{enumerate}
\end{lemma}
\begin{proof}
A node at depth $q$ occurs $T^q$ times in the root word, each time in
a word of length $(mT)^{k-q}$. Its occurrences occupy the fraction
$m^{-q}$ of all positions. Every occurrence receives its subtree phase
vector plus fixed offsets from its ancestors, and those offsets are
the same on both sides of each comparison.

Within an occurrence of $X_v$, addition of $h_v(a)$ changes the row
indices from $0,\ldots,T-1$ to $a,\ldots,T-1+a$. Using the representative
$0\leq a<P_v\leq k^{48k}<T$, the lists share $T-a$ rows. Deleting $a$
rows and inserting $a$ rows costs at most $2a/T$ times the occurrence
length. Nothing changes outside the occurrences of $v$. This proves
\eqref{eq:cycle-cost}; it requires no divisibility of $T$ by $P_v$.

For a tiny rotation, each affected symbol quantizes a uniformly
distributed leaf phase plus a fixed offset. A rotation by
$0\leq\eta\leq1/Q$ changes its interval with probability
$R\eta\leq R/Q$. Sum these positional mismatch probabilities and use
substitutions to prove \eqref{eq:tiny-cost}. Linearity of expectation
suffices even when different occurrences use the same phase. This
estimate is an average over $w$. Finally, an arbitrary change of one
leaf can affect only its $Nm^{-k}$ occurrences, proving
\eqref{eq:leaf-cost} pointwise.
\end{proof}

\subsection{The Fourier identity on the torus}

The next identity expresses averaged $\ellone$ distances through
nonnegative frequency weights. It is the continuous counterpart of
Lemma~\ref{lem:finite-fourier}; its averaging measure and its integer
frequencies are different.

\begin{lemma}[Torus displacement]\label{lem:torus-fourier}
Let $n\geq1$ be an integer and let $F:\T^n\to\ellone$ be a measurable
map with finite range. Set
\[
 L_F(h)=\E_w\|F(w+h)-F(w)\|_1
 \qquad(h\in\T^n),
\]
using normalized Haar measure. There are summable coefficients
$\mu_\xi\geq0$, indexed by $\xi\in\Z^n$, such that
\begin{equation}\label{eq:torus-fourier}
 L_F(h)=\sum_{\xi\in\Z^n}
 \mu_\xi\bigl(1-\cos(2\pi\langle\xi,h\rangle)\bigr)
 \qquad(h\in\T^n).
\end{equation}
\end{lemma}
\begin{proof}
Apply Lemma~\ref{lem:cuts} to the finite range of $F$ and pull a cut
back to a measurable indicator $U$ on $\T^n$. It belongs to
$L^2(\T^n)$. For
$\widehat U(\xi)=\int_{\T^n}U(w)e^{-2\pi i\langle\xi,w\rangle}\,dw$,
Parseval and the Fourier translation rule give
\[
 \int_{\T^n}|U(w+h)-U(w)|^2\,dw
 =\sum_{\xi\in\Z^n}2|\widehat U(\xi)|^2
 \bigl(1-\cos(2\pi\langle\xi,h\rangle)\bigr).
\]
An indicator difference has absolute value equal to its squared
absolute value. Multiply the identity by each cut weight and sum over
the finitely many cut types. The resulting coefficients are nonnegative.
Their sum is at most twice the total cut weight, because
$\sum_\xi|\widehat U(\xi)|^2=\int|U|^2\leq1$. This also proves
summability and justifies the identity for every $h$.
\end{proof}

\subsection{Two tests and forced branching}

Apply the notation of Lemma~\ref{lem:torus-fourier} with $n=m^k$.
For a frequency $\xi\in\Z^{m^k}$ and a node $v$, define its integer
summary
\[
 s_v=\langle\xi,e_v\rangle
 =\sum_{u\text{ leaf below }v}\xi_u.
\]
Thus $s_v=\sum_i s_{vi}$ at an internal node and $s_u=\xi_u$ at a
leaf. The multiplier of the root half-turn is $1-\cos(\pi s_o)$,
equal to two for odd $s_o$ and zero for even $s_o$. We bound it using
tests that cover every frequency with $s_o\ne0$.

\begin{lemma}[Torus detection inequality]\label{lem:detection}
For the fixed construction and its shifts $h_v(a)$, every measurable
finite-range $F:\T^{m^k}\to\ellone$ satisfies
\begin{align}
 L_F(e_o/2)\leq{}&
 4\sum_{v\text{ all nodes}}
 \E_{\eta\in[0,1/Q]}L_F(\eta e_v)
 +2\sum_{v\text{ internal}}
 \E_{a\in\Z/P_v\Z}L_F(h_v(a))\notag\\
 &\quad+2k^{-k}\sum_{u\text{ leaf}}
 \E_{\eta\in\T}L_F(\eta e_u).
 \label{eq:detection}
\end{align}
Every interval, finite-group and circle average has total mass one.
\end{lemma}
\begin{proof}
We compare the multipliers of a fixed frequency in
\eqref{eq:torus-fourier}. If $|s_v|\geq Q$ at some node, then
\[
 Q\int_0^{1/Q}(1-\cos(2\pi s_v\eta))\,d\eta
 =1-\frac{Q}{2\pi s_v}\sin(2\pi s_v/Q)
 \geq1-\frac1{2\pi}\geq\frac12.
\]
The first sum in \eqref{eq:detection}, with coefficient four, therefore
dominates the root multiplier, which is at most two.

If an internal node satisfies
\begin{equation}\label{eq:nonzero-congruence}
 \sum_i c_{vi}s_{vi}\not\equiv0\pmod{P_v},
\end{equation}
then the average of the multiplier for $h_v(a)$ over
$a\in\Z/P_v\Z$ is exactly one, by the sum of a nontrivial cyclic
character. The second sum, with coefficient two, dominates the root
multiplier in this case.

It remains to consider a frequency with $|s_v|<Q$ at every node and
with every internal congruence equal to zero. If $s_o=0$, its root
multiplier vanishes. Suppose $s_o\ne0$, and consider an internal node
$v$ with $s_v\ne0$. Call a group of children empty if all summaries
in that group vanish. For an empty group $j$, reduce the zero
congruence modulo $r_{v,j}$. The coefficients outside that group are
one modulo $r_{v,j}$ and the group contributes zero, so
$r_{v,j}\mid s_v$.

If at least $g/2=8k$ groups were empty, their distinct primes would
have product at least
$(k^2)^{8k}=k^{16k}>Q$. That product cannot divide a nonzero integer
$s_v$ with $|s_v|<Q$. Hence more than $g/2$ groups contain a child
with nonzero summary, and in particular there are at least $k$ such
children. Starting from the root and repeating the conclusion gives
at least $k^k$ leaves $u$ with $s_u\ne0$. At each such leaf, the
multiplier for $\eta e_u$ has mean one over uniform $\eta\in\T$.
The third sum in \eqref{eq:detection}, with coefficient $2k^{-k}$,
is therefore at least two.

These cases bound the root multiplier for every frequency. Multiply
by the nonnegative coefficients $\mu_\xi$ and sum. Their summability
justifies interchange with all normalized shift averages, proving
\eqref{eq:detection}.
\end{proof}

The cutoff and congruence tests have different jobs. The divisibility
argument forces branching only after all summaries are known to have
absolute value less than $Q$. It would not rule out a large nonzero
integer divisible by many of the selected primes.

\subsection{The binary obstruction and every length cap}

Apply Lemma~\ref{lem:delimiter} to the torus alphabet of size
$A=m^kR$, choosing labels of length $\lceil\log_2A\rceil$. Write
$\mathcal B$ for the resulting encoding. Its block width is
\begin{equation}\label{eq:torus-code-width}
 B=2\lceil\log_2(m^kR)\rceil+3=O(k\log k).
\end{equation}
The encoded image consists of binary words of one length $BN$.
It has at least two points, since the lower code bound preserves the
positive root separation in \eqref{eq:torus-separation}. We apply the
detection inequality directly to an embedding of this encoded image;
the code's width will enter only the upper cost estimates.

Take any embedding $f$ of the encoded image with distortion $D$ and
rescale it to be noncontracting with expansion at most $D$. Set
\[
 F(w)=f(\mathcal B(X(w))),\qquad L(h)=L_F(h).
\]
The map $F$ is measurable with finite range, so
Lemma~\ref{lem:detection} applies. The pointwise separation, the three
costs in Lemma~\ref{lem:torus-costs}, and the code bounds give
\begin{align}
 L(e_o/2)&\geq\gamma^kN/2,\label{eq:torus-far}\\
 L(h_v(a))&\leq2DBNm^{-q}P_v/T
 &&(v\text{ internal},\ a\in\Z/P_v\Z),\label{eq:torus-cycle-l1}\\
 L(\eta e_v)&\leq DBNm^{-q}R/Q
 &&(v\text{ any node},\ 0\leq\eta\leq1/Q),\label{eq:torus-tiny-l1}\\
 L(\eta e_u)&\leq DBNm^{-k}
 &&(u\text{ leaf},\ \eta\in\T).\label{eq:torus-leaf-l1}
\end{align}
For the second line, use the representative $0\leq a<P_v$ in the
pointwise cost proof; the displacement itself depends only on the
residue. The third line uses the Haar-average cost, with no pointwise
claim about tiny rotations.

There are $m^q$ nodes at depth $q$. Thus the factors $m^{-q}$ cancel
when the upper bounds are summed over a depth in
\eqref{eq:detection}. We obtain
\begin{align*}
 L(e_o/2)
 &\leq DBN\left(
 4(k+1)\frac RQ+4k\frac{k^{48k}}T+2k^{-k}\right)\\
 &\leq C'DBNk^{-k},
\end{align*}
for an absolute $C'$, because $R/Q=k^{-4k}$ and
$k^{48k}/T=k^{-4k}$. Comparing this with \eqref{eq:torus-far} gives
\begin{equation}\label{eq:torus-final-D}
 D\geq\frac{(\gamma k)^k}{2C'B},
 \qquad\gamma=1/128.
\end{equation}
The denominator retains the width $B$ from the delimiter code.
Since $B=O(k\log k)$, the logarithm of the right side is at least
$c_{\mathrm{tor}}k\log k$ for an absolute $c_{\mathrm{tor}}>0$ and all
sufficiently large $k$.

The binary word length satisfies
\begin{equation}\label{eq:torus-log-length}
 \log(BN)=52k^2\log k+O(k\log k).
\end{equation}
For every sufficiently large integer $d$, put $\tau=\log d$ and choose
$k=\lfloor a\sqrt{\tau/\log\tau}\rfloor$, where $a>0$ is a sufficiently
small absolute constant. There is an absolute $C_2$ such that
$\log(BN)\leq C_2k^2\log k$ for all large $k$. Since
$k\leq a\sqrt{\tau/\log\tau}$ and $\log k\leq\log\tau$ eventually, choosing
$C_2a^2\leq1/2$ gives $\log(BN)\leq\tau/2<\tau$.
Also, for all large $\tau$,
\[
 k\geq\frac a2\sqrt{\frac \tau{\log\tau}},
 \qquad \log k\geq\frac13\log\tau.
\]
It follows that $k\log k\geq(a/6)\sqrt{\tau\log\tau}$.
The finite set
$\{\mathcal B(X(w)):w\in\T^{m^k}\}\subseteq\{0,1\}^{BN}$ therefore
has distortion at least
$\exp(c\sqrt{\log d\,\log\log d})$ and has $BN\leq d$.
This completes the torus proof of Theorem~\ref{thm:binary} for every
sufficiently large cap.
